\documentclass[12pt]{article}
\usepackage[utf8]{inputenc}
\usepackage[T1]{fontenc}
\usepackage{hyperref}
\title{Wandarang Language Resources for K-12 Curriculum Development}
\author{Compiled by Claude Sonnet 5.5 (claude-sonnet-5-5) for Andr\'as Kornai}
\date{}
\begin{document}
\maketitle

\section{Introduction}
Wandarang (also spelled Warndarang or Warndarrang; ISO 639-3 wnd) is a non-Pama-Nyungan language of the Northern Territory, Australia. It is catalogued by AIATSIS as N120. It is generally regarded as extinct. One search summary said it was thought extinct by 1974 but that a fluent speaker was reportedly recorded in 1989, giving oral texts with a Kriol translation. I did not verify that claim at its source. Its classification is disputed: some sources link it to Marra and Alawa in a Marran grouping, while Wikipedia places it in Gunwinyguan.

\section{Primary Educational Resources}
I found no primer, textbook, school programme or classroom material for Wandarang. A heritage curriculum would have to be built from archival and linguistic sources, preferably with descendant communities.

\section{Linguistic and Reference Materials}
The AIATSIS Austlang entry gives basic identification and links to archival holdings: \url{https://collection.aiatsis.gov.au/austlang/language/N120}. An encyclopedic overview is at \url{https://wikipedia.com/wiki/Warndarrang}. The AfBo survey of affix borrowing lists Warndarrang, which is useful for contact linguistics: \url{https://afbo.info/languages/warndarangwnd}. I did not find a published grammar or dictionary in the search results.

\section{Religious and Community Resources}
I found no scripture translations or community organisation pages specific to Wandarang.

\section{Digital and Audio Resources}
No audio or digital collection was seen directly in the search results beyond the AIATSIS catalogue entry above. Recordings may exist in archives, but I did not confirm any.

\section{Curriculum Development Context}
Very little is documented. Any curriculum should be a revitalization or awareness programme: place names, the history of the language, and archival materials, developed in consultation with the relevant Aboriginal communities. The first step would be to ask AIATSIS about what its holdings contain and what access conditions apply.

\end{document}